\section{Related Work}
\label{sec:related}
\subsection{Other datasets used in neural abstractive summarization}
Neural abstractive summarization was pioneered in
\citet{Rush2015}, where they train headline generation models using the
English Gigaword corpus
\citep{graff2003english}, consisting of news articles from number of publishers.
However, the task is more akin to sentence paraphrasing
 than summarization
as only the first sentence of an article is used to predict the headline,
another sentence.
RNN-based encoder-decoder models with attention (seq2seq) perform very well on this task
in both
ROUGE \citep{lin2004rouge}, an automatic metric often used in summarization,
and human evaluation
\citep{chopra2016abstractive}.

In \citet{nallapati2016abstractive}, an abstractive
summarization dataset is
proposed by modifying a question-answering dataset of news articles paired
with story highlights from Daily Mail and CNN. This task is more difficult than
 headline-generation because the information used in the highlights may
 come from many parts of the article and not only the first sentence. One downside
 of the dataset is that it has an order-of-magnitude fewer parallel examples (310k vs. 3.8M) to learn from.
 Standard seq2seq models with attention do less well, and a number
 of techniques are used to augment performance.
 Another downside is that it is unclear what the guidelines are for
 creating story highlights and it is obvious that there are significant
 stylistic differences between the two news publishers.

 In our work we also train neural abstractive models, but
 in the multi-document regime with Wikipedia.
 As can be seen in Table~\ref{dataset-sizes}, the input and output text are
 generally much larger, with significant variance depending on the article.
 The summaries (Wikipedia lead) are multiple sentences and sometimes multiple
 paragraphs, written in a fairly uniform style as encouraged
 by the Wikipedia Manual of Style\footnote{\url{https://en.wikipedia.org/wiki/Wikipedia:Manual_of_Style}}. However, the input documents may consist
 of documents of arbitrary style originating from arbitrary sources.

 We also show in Table~\ref{dataset-sizes} the ROUGE-1 recall scores of the output
 given the input, which is the proportion of unigrams/words
 in the output co-occuring
 in the input. A higher score corresponds to a dataset more amenable to
 extractive summarization. In particular, if the output is completely embedded somewhere
 in the input (e.g. a wiki-clone), the score would be 100. Given a score of only 59.2 compared
 to 76.1 and 78.7 for other summarization datasets shows that ours is the least amenable to purely extractive methods.

\subsection{Tasks involving Wikipedia}
There is a rich body of work incorporating Wikipedia for machine learning tasks, including
question-answering (\citet{hewlett2016wikireading}, \citet{rajpurkar2016squad}) and
information extraction \citep{lehmann2015dbpedia}, and text generation from 
structured data 
\citep{Lebret2016}.

The closest work to ours involving generating Wikipedia is \cite{Sauper2009},
where articles are generated extractively (instead of abstractively in our case)
from reference documents using learned
templates. The Wikipedia articles are restricted to two categories, whereas we
use all article types. The reference documents are obtained from a search engine,
with the Wikipedia topic used as query similar to our search engine references.
However we also show results with documents only found in the
References section of the Wikipedia articles.

\subsection{Transformer models} \label{related_transformer}
Previous work on neural abstractive summarization relies on RNNs as
fundamental modules, mirroring techniques successful
in machine translation (MT).
Recently, state-of-the-art MT results were obtained using a
non-recurrent architecture, called the Transformer \citep{vaswani2017attention}.
The lack of recurrence enables greater within-training-example parallelization,
at the cost of quadratic complexity in the input sequence length.
We find the Transformer transfers well to medium length, input sequence
summarization and describe modifications to better handle longer sequences.

\begin{table}[t]
\caption{Order of magnitude input/output sizes  and
	unigram recall for summarization datasets.}
\label{dataset-sizes}
\begin{center}
\begin{tabular}{lllll}
\multicolumn{1}{c}{\bf Dataset}  &\multicolumn{1}{c}{\bf Input}
	&\multicolumn{1}{c}{\bf Output}
	&\multicolumn{1}{c}{\bf \# examples}
	&\multicolumn{1}{c}{\bf ROUGE-1 R}
\\ \hline \\
	Gigaword \citep{graff2003english}         & $10^1$         & $10^1$ & $10^6$ & $78.7$\\
	CNN/DailyMail \citep{nallapati2016abstractive} & $10^2$--$10^3$          & $10^1$  & $10^5$ & $76.1$\\
	WikiSum (ours)          & $10^2$--$10^6$  & $10^1$--$10^3$ & $10^6$ & $59.2$\\
\end{tabular}
\end{center}
\end{table}

